\section{Questions}

Much of the program's worth lies in what it leaves open, stated plainly rather than buried.

The sharpest is the self on the bare substrate. The capture, the bath-made self, and the nesting were all shown cleanly in a reduced model, but on the committed lattice a single local rule can confine a body or radiate from it, not both at once, so a self native to the bare five base things, rather than to a reduced caricature, is not yet in hand. Closely tied to it is a precise missing ingredient in the field theory. The bare single-charge rule yields only diffusive modes, spreading as the square root of time, while a relativistic massless particle, a photon or a graviton or a wave of sound, needs a sharp light-cone mode, and getting one appears to require a second conserved quantity, a momentum or inertia, beyond the single charge. Whether the base supplies it is open.

Several numbers and laws sit just past reach. The three generations are forced as a threefold structure but their three different masses are not, the boldest open bet. The overall strength of the coupling, the fine-structure constant, is left free by the bare law. Gravity in the curved four-dimensional bulk, as opposed to the flat cusp where Newton and Einstein are clean, is not yet derived, and even on the cusp the attraction that does the work now is an effective, instantaneous one, so a stable, reversible, propagating gravity built at the base is a gap of its own. Dark matter has a candidate geometric form, the bulk's nonlocal pull, but not yet a sharp prediction to test, and the discrete symmetries, present and consistent to fifteen digits, still want a formal proof of their combined invariance. The growth of the committed honeycomb is measured to a rate near eighteen but has no closed-form recurrence. And the full reconstruction of quantum mechanics, the Bell bound reached and the Born rule forced, still wants a derivation of measurement and collapse rather than an assumption of them.

Two questions are deeper still. Whether a pattern written into the bulk truly lasts against the churn, the persistence problem, is the one frontier on which the realms and beings and the survival of anything after death all quietly depend. And whether the world began or is eternal, both consistent with the law, is left genuinely undecided, since the necessity of a beginning, argued earlier, fixes that there was a first state without fixing whether the heap of distinctions runs back forever or starts from the bootstrap. The model points past where it can yet walk, and the pointing is marked as such.

One question runs underneath all of these, whether the base itself is forced. The dock is already a convergence, the same twenty-four points reached as a shell of the tone alphabet, as a root system, as a group of quaternions, and as the tightest sphere packing, four roads to one figure, and the dimension four is the only one where spinors, self-duality, and triality arrive together. So the base may not be a list of choices at all but the single unavoidable shape of a world made to feel itself, the ternary tone, the four directions it spans, the 24-cell they shell into, and the one reversible law that conserves them, each following from the last. Whether the whole substrate can be derived this way from a single seed, made a theorem rather than a posit, is the deepest question the program holds, and the one whose answer would change what kind of thing the theory is.
